\begin{abstract}
\section{}
Handcrafted leaf measurements are easy to inspect, but it is unclear what a graph built over them adds to stress recognition. We examined this question with the Open Leaf Image Dataset I (OLID-I; 4,749 photographs of eight vegetable crops) using a binary endpoint: healthy versus any recorded stress. An excess-green segmentation accepted 4,387 photographs, from which we extracted 38 color, shape, texture and regional descriptors. We compared Weisfeiler--Lehman (WL) subtree kernels on five descriptor graphs with graph-free kernels, random graphs, graph neural networks, descriptor classifiers and frozen ImageNet encoders on shared folds. Three refinements on a hybrid graph raised pooled macro-F1 from 0.596 (descriptor states alone, tuned SVM cost) to 0.697. Tuning the SVM cost did not explain this gain. At depth three, 96\% of training patterns occurred in only one image, and a control reproducing that effect without any shared pattern recovered about a third of the gain (0.629). A kernel counting matched triples of descriptor states, with no graph, reached 0.690. Across 200 random graphs with the hybrid edge count, macro-F1 ranged from 0.680 to 0.712, and 88 of them matched or exceeded the hybrid graph. Under nested selection, WL (0.685 with crop supplied, 0.682 with predicted crop) did not exceed the graph-free triple kernel (0.699, 0.683) or tuned Extra Trees (0.742, 0.729); frozen EfficientNet-B0 led with predicted crop (0.739). Segmentation determined which leaves could be screened: 100 of 207 leaf-miner photographs were rejected, so the descriptor pipeline gave a positive screen for 42.5\% of them, compared with 90.8\% for a frozen encoder trained on the same photographs but scoring every submission. Across simulated acquisition changes in all five test folds, descriptor models lost 0.08--0.15 macro-F1 on average and WL was not consistently the most robust of them ($-0.088$, nearest neighbors $-0.083$, tuned Extra Trees $-0.095$), whereas frozen encoders lost almost nothing ($-0.003$ for EfficientNet-B0). The measured benefit of the graph kernel therefore lies mainly in jointly matching discretized descriptors rather than in the proposed links between them. Screening studies should also report how many submitted photographs reach the classifier.

\tiny\keyFont{\section{Keywords:} plant phenotyping, leaf images, plant stress, graph kernels, Weisfeiler--Lehman, segmentation, benchmark, interpretability}
\end{abstract}

\section{Introduction}
Image-based phenotyping promises fast, repeatable screening of plant stress, but the path from photograph to decision contains several choices that are rarely evaluated separately \citep{Fahlgren2015,Minervini2015,Singh2016}. Deep convolutional networks now dominate image-based disease recognition \citep{Mohanty2016,Ferentinos2018,Kamilaris2018}, yet their accuracy depends on dataset size and variety and degrades under acquisition conditions not represented in training \citep{Barbedo2016,Barbedo2018}. Pipelines built on explicit measurements, such as those produced by PlantCV \citep{Gehan2017}, remain attractive because each number has a physical meaning: if mean hue falls, a researcher can check whether the change came from the leaf center or margin. Methods that explain network decisions exist \citep{Ghosal2018}, but a measurement pipeline makes its intermediate quantities directly available for scrutiny \citep{Lobet2017}.

A measurement pipeline also introduces a step that pixel classifiers can avoid: the leaf must first be segmented. Color indices such as excess green separate vegetation from background well when tissue is green \citep{Woebbecke1995,Meyer2008}, but stressed tissue is often yellow, brown or perforated. A survey can therefore miss damage in two ways: segmentation can reject the photograph, or the classifier can label an accepted leaf as healthy. Evaluating only accepted photographs counts the second kind of miss.

Graph kernels offer one way to compare several measurements jointly \citep{Kriege2020,Borgwardt2020}. \citet{Sriram2026} encoded image-derived plant traits as a graph for plant-condition recognition. In such a design each measurement becomes a node with a discretized label, edges connect related measurements, and the Weisfeiler--Lehman (WL) subtree kernel compares leaves by the labeled neighborhoods they share \citep{Shervashidze2011}. Two leaves with similar marginal color states could still be distinguished if a color state co-occurred with a particular texture state. Graph neural networks (GNNs) generalize the same neighborhood aggregation with learned functions and are at most as discriminative as the WL test \citep{Xu2019,Morris2019}.

The central uncertainty is whether the proposed edges contribute anything beyond the generic effect of comparing several discretized descriptors at once. A graph kernel can also gain or lose performance for reasons unrelated to the graph, for example when most refined patterns occur only once and the kernel matrix becomes dominated by self-similarity \citep{Yanardag2015,Greene2006}. We therefore asked four questions using OLID-I, an expert-annotated collection of eight vegetable crops with healthy, nutritional, pest, disease and combined conditions \citep{Orka2023}:
\begin{enumerate}
\item What does neighborhood refinement add, once the SVM cost, the scale effect of unique patterns, graph-free joint matching and random graphs with identical edge counts are controlled?
\item How does the WL kernel compare with descriptor classifiers, GNNs and frozen image encoders when all models are tuned in the same nested way?
\item Which stress conditions are lost at segmentation, and how much of that loss does a classifier that needs no mask recover?
\item How do simulated acquisition changes affect both the acceptance and the classification of photographs?
\end{enumerate}

\section{Materials and methods}
\subsection{Image collection, labels and coverage}
OLID-I contains 4,749 RGB images of vegetable leaves sampled at the Bangladesh Agricultural Research Institute in Gazipur. The dataset authors photographed individual leaves on black fabric with an iPhone 13 Pro Max at $3024\times3024$ pixels. They held the camera distance constant within each crop and specified different distances across crops \citep{Orka2023}. The eight crops were ash gourd, bitter gourd, bottle gourd, cucumber, eggplant, ridge gourd, snake gourd and tomato. We mapped the healthy label to zero and each of the other recorded conditions, including combined conditions, to one. The resulting endpoint indicates whether the source annotation records any stress. It does not quantify severity or identify its cause.

For normalized red, green and blue intensities $r,g,b$, the segmentation score was the excess-green index $2g-r-b$ \citep{Woebbecke1995}. We rescaled each image's score to eight bits, applied Otsu thresholding \citep{Otsu1979}, performed opening and closing with a $5\times5$ elliptical element (two iterations each), and filled the largest external contour. The accepted interval for mask area divided by image area was $[0.01,0.98]$. This check detects extreme mask sizes; it cannot establish boundary accuracy. Because the index rewards green tissue, yellow, brown or perforated lamina is the expected failure mode. We calculated retention from all source records. All descriptor models used the accepted cohort; image encoders were evaluated on both the accepted cohort and every submitted photograph (Section~\ref{sec:coverage-methods}).

Figure~\ref{fig:leaf} shows three photographs of the same crop with their computed masks. OLID-I already provides a controlled dark background; introducing a segmentation stage is a choice of this measurement pipeline \citep{Orka2023}.
\showpanel{\leafpanel}

\subsection{Descriptor extraction}
Ten measurements characterized contour geometry, including area, circularity, solidity and fitted ellipse axes. Eighteen summarized the means and standard deviations of RGB, HSV and CIELAB components within the mask. Two ratios summarized green and yellow color. Four gray-level co-occurrence statistics summarized contrast, homogeneity, energy and correlation \citep{Haralick1973}. These used one-pixel offsets at four angles in the bounding rectangle, which includes black background around the leaf; the texture statistics therefore partly reflect the leaf outline. The last four measurements were mean hue and a yellow ratio for each of two regions, the masked central third in both dimensions and the remaining leaf area. Supplementary Data Sheet S7 defines the individual descriptors. Pixel lengths and areas retained the acquisition scale, so they also carry the crop-specific camera distance.

\subsection{Descriptor graphs and the WL kernel}
Each graph contained 38 named nodes. A node represented the same descriptor in every image. Training-fold quantiles encoded its value as low, middle or high. The adjacency was shared by all images processed by a fitted model and was learned from training images only. We compared five undirected constructions. The correlation graph joined pairs with absolute training correlation of at least 0.70. The nearest-neighbor graph standardized each descriptor across training images and joined it to its three closest descriptor vectors by Euclidean distance, symmetrizing the edges. The family graph connected all members of each descriptor family and added the strongest absolute-correlation link between every pair of families. The regional graph connected the two core descriptors, the two edge descriptors and corresponding core--edge pairs, then linked other nodes to regional features at absolute correlation of at least 0.60. The hybrid graph combined the family graph with correlation edges at 0.80 and regional edges at 0.70. Figure~\ref{fig:matrices} shows the matrices fitted on the first training fold. An edge is an operation on measured variables; it has no spatial coordinates on the lamina, and RGB, HSV and Lab summarize overlapping color information.
\showpanel{\matrixpanel}

At iteration zero, node labels encoded descriptor identity and quantile. If $\ell_v^{(t)}$ is the current label of descriptor $v$ and $N(v)$ its neighbors, the updated label is an index for the signature
\begin{equation}
  \left(\ell_v^{(t)},\,\operatorname{sort}\{\ell_u^{(t)}:u\in N(v)\}\right).
\end{equation}
Let $\phi_h(G)$ concatenate counts of the initial labels and refinements through depth $h$. The kernel was
\begin{equation}
 K_h(G_i,G_j)=\frac{\phi_h(G_i)^\mathsf{T}\phi_h(G_j)}
 {\|\phi_h(G_i)\|_2\,\|\phi_h(G_j)\|_2}.
\end{equation}
Bins, data-dependent edges and the pattern vocabulary were learned on the training partition. Refinement patterns absent from that vocabulary received no coordinate. A class-balanced support vector machine (SVM) used the precomputed kernel \citep{Cortes1995}. At $h=0$ the kernel compares discretized measurements directly, so varying $h$ from zero to three isolates the added neighborhood information. The original study fixed the SVM cost at $C=1$; here we also tuned $C\in\{0.01,0.1,1,10,100\}$.

\subsection{Controls for the refinement gain}
\label{sec:controls}
Four controls separated components of the refinement gain. First, the SVM cost was tuned at every depth, so that a gain could not arise from a cost suited to one kernel but not another. Second, we recorded for every held-out image which refined labels existed in the training vocabulary, and for every training label how many training images carried it.

Third, a \emph{private-pattern control} reproduced the scale effect of unique deep patterns without any shared neighborhood. If no refined pattern recurs, each training graph has $38$ initial labels plus up to $3\times38$ refined labels found in no other image, while a held-out graph has only its $38$ initial labels in the vocabulary. With $m_{ij}$ the number of descriptors in the same quantile state for images $i$ and $j$, the corresponding kernel is $m_{ij}/152$ between distinct training images, one on the training diagonal, and $m_{tj}/\sqrt{38\times152}$ between a held-out image $t$ and a training image. Any gain of this control over depth zero arises from the changed balance between self-similarity and shared states, not from neighborhood information.

Fourth, two graph-free \emph{joint-matching kernels} compared states without any edges. With $x_i$ the one-hot vector of the 38 quantile states, $K(i,j)=(x_i^\mathsf{T}x_j/38)^d$ counts matched pairs ($d=2$) or triples ($d=3$) of descriptor states, normalized by the number of descriptors. These kernels test whether comparing combinations of states, without choosing which descriptors to combine, explains the WL gain.

Finally, three random-graph null models replaced the hybrid adjacency fitted in each training fold: (i) permuted descriptor positions, $PAP^\mathsf{T}$ for a random permutation $P$, which keeps the graph's shape but reassigns which measurement occupies each position; (ii) degree-preserving rewiring by ten successful double-edge swaps per edge, which keeps each descriptor's degree but changes its neighbors; and (iii) a random graph with the same number of edges. We fitted 100 permuted, 50 rewired and 50 random graphs, each on all five outer folds at depth three and $C=1$ (the value selected for the hybrid graph by nested tuning). Descriptor values, bins, folds and labels were unchanged. We report each graph's five-fold mean and the share of null graphs scoring at least as well as the hybrid graph.

\subsection{Prediction tasks and comparators}
The pooled task used one binary classifier for all crops. The known-crop task trained and evaluated a separate classifier within each crop. The routed task first predicted a crop and applied the binary model fitted to training leaves of that crop; a parallel evaluation supplied the true crop to the same classifiers to isolate the cost of routing errors. Descriptor pipelines used a 300-tree random forest for routing \citep{Breiman2001}. For the image pipelines, a logistic classifier routed the frozen embedding and a crop-specific logistic classifier predicted stress.

Descriptor comparators included linear and radial-basis SVMs, logistic regression, nearest neighbors, random forest, Extra Trees, histogram gradient boosting, XGBoost \citep{Chen2016}, LightGBM, AdaBoost, a decision tree, Gaussian naive Bayes and a multilayer perceptron (MLP). The image comparators used ImageNet-pretrained MobileNetV3-Large \citep{Howard2019} and EfficientNet-B0 \citep{Tan2019} as frozen feature extractors \citep{Kornblith2019}. Images were decoded at reduced size, resized to $224\times224$, converted to RGB and normalized by the ImageNet channel constants. The encoders supplied 960 and 1,280 pooled features, respectively. Standardization and class-balanced logistic regression were fitted to training embeddings, with the regularization parameter selected from $\{0.001,0.01,0.1\}$ by inner cross-validation.

Pooled and routed tasks used the same five stored image partitions stratified by crop--condition class. Known-crop evaluation used shuffled binary-stratified five-fold splits within each crop (seed 42). Image-level test folds were held out from all fitting. Plant and acquisition-session identifiers were unavailable for grouped evaluation.

\subsection{Model selection}
We distinguish two WL estimates. The \emph{selected} estimate fixes the graph and depth chosen in an earlier sweep on all study folds (hybrid, depth three, for the pooled task; one setting per crop otherwise). It measures a retrospectively chosen configuration and is optimistic \citep{Varma2006,Cawley2010}. The \emph{nested} estimate selects the configuration by macro-F1 in three stratified inner folds of each outer training partition. The original nested grid crossed five graphs with depths one to three at $C=1$; the extended grid added depth zero and the five cost values. We ran nested selection for all three tasks, the graph-free kernels and the private-pattern control. The rerun reproduced the archived known-crop nested WL scores exactly on all 40 crop--fold splits. For the routed task the router was refitted; it assigned the correct crop to 91.1\% of held-out images, compared with 91.9\% in the archived run. Five continuous-feature comparators were also tuned in inner folds; for example, the 300-tree Extra Trees grid crossed two feature-sampling settings (square root or one half) with minimum leaf sizes of one, three or five. The nested XGBoost grid covered 200 or 500 trees, depth three or six, learning rate 0.05 or 0.10 and subsampling of 0.8 or 1.0. Because the selected WL estimate is optimistic, we base comparisons between methods on nested estimates.

\subsection{Learned graph baselines and tree attributions}
We compared three GNNs on the same 38 descriptor nodes and outer image partitions: a graph convolutional network (GCN) with degree-normalized aggregation \citep{Kipf2017}, a graph isomorphism network (GIN) that sums neighbors and applies a multilayer perceptron \citep{Xu2019}, and a graph attention network (GAT) with four attention heads \citep{Velickovic2018}. Each node received its standardized measurement clipped to $[-5,5]$, a one-hot quantile bin and a learned eight-dimensional descriptor-identity embedding. The networks used three layers of width 32 with residual connections, layer normalization and dropout of 0.2, and combined mean and maximum pooling over nodes from the input projection and each layer. Training used Adam (learning rate $10^{-3}$, weight decay $10^{-4}$), batches of 64 and class-weighted cross-entropy for 20--100 epochs with patience 20; a stratified 10\% subset of the training data selected the checkpoint. In the smallest crops this validation subset contained only a few healthy leaves, so checkpoint selection was noisy. Nested GIN searched five graphs and two or three layers under the same inner protocol as nested WL; GCN and GAT were evaluated at default settings and should be read as secondary baselines. A pooled graph study refitted WL (depth three) and the networks (three layers) on each graph, and for the networks a node-only control removed inter-node edges while keeping own-node operations and readout. We computed TreeSHAP contributions from default pooled XGBoost on each outer test fold \citep{Lundberg2020}.

\subsection{Coverage with an encoder that needs no mask}
\label{sec:coverage-methods}
To measure what segmentation costs, we embedded all 4,749 photographs with both frozen encoders. Accepted photographs kept their stored folds. Each of the 362 rejected photographs was assigned to an outer fold by shuffling within its crop--condition class (seed 42) and dealing folds in turn, giving 72--73 rejected photographs per fold. Encoder classifiers were trained on the accepted training photographs only, exactly as in the main comparison, and then scored every photograph of the held-out fold; only coverage therefore changed. A second variant also trained on rejected training photographs. For each stress condition with $n$ original photographs, $a$ accepted photographs and $t$ positive calls, $t/a$ measures sensitivity among accepted inputs, whereas $t/n$ measures the share of submitted photographs that yielded a positive screen. A rejected photograph counts as no call for the descriptor pipeline. Descriptor-pipeline calls came from the pooled hybrid model at depth three, so decisions had the same meaning across crops.

\subsection{Metrics, statistics and perturbations}
Macro-F1, the mean of the healthy and stressed F1 scores, was the primary endpoint. We also report balanced accuracy, healthy recall (specificity), stressed recall (sensitivity) and accuracy, together with the always-stressed reference. Pooled and routed results are means and standard deviations of outer-fold scores. Known-crop scores first average folds within crop and then give each crop equal weight. For matched crop--fold scores we report mean differences, fold wins and two-sided Wilcoxon signed-rank tests on the 40 fold differences and, because folds within a crop share training images, exact tests on the eight crop-mean differences \citep{Dietterich1998,Nadeau2003}. All $p$-values are unadjusted and exploratory.

To characterize the measurements, we calculated a signed rank contrast for each crop and descriptor: twice the single-descriptor ROC area minus one, with ties receiving half credit. It equals the probability that a randomly drawn stressed image from the crop has the larger value minus the probability that a healthy one does. This descriptive analysis used all accepted leaves and never informed model selection.

Learning curves used stratified fractions of the outer training partition (10, 20, 35, 50, 75 and 100\%) with the original held-out tests, at fixed settings. Processing times combine measured image preparation and single-image inference; descriptor extraction used full-resolution images, whereas the CNN path used reduced JPEG decoding.

The robustness experiment sampled 240 photographs from each of the five pooled test folds (1,200 in total), stratified by crop and stress status; the fold-zero sample reproduced the 240 photographs of the original single-fold study. Models trained on clean images of the other four folds. Each test image was independently darkened or brightened by 25\%, given a warm cast (red multiplied by 1.12 and blue by 0.85), Gaussian-blurred ($\sigma=4$ pixels), JPEG-compressed at quality 20, perturbed with Gaussian noise ($\sigma=12$ intensity units) or halved in both dimensions. Each change preceded segmentation and measurement, or encoder resizing. For each model and fold we averaged the seven differences from clean macro-F1, and we summarize these across folds. Descriptor pipelines scored only images that passed segmentation after each change; we therefore also recomputed differences on images with predictions from every model under every change. The five-fold study covered the WL kernel, eleven descriptor classifiers and both encoders; the GNN robustness results come from the original single-fold study.

\section{Results}
\subsection{Coverage of leaves and stress conditions}
The filter accepted 4,387 images (92.4\%) and excluded 362. Among accepted images, 3,674 were labeled stressed and 713 healthy. Of 786 healthy images in the original collection, 713 passed (90.7\%), compared with 3,674 of 3,963 stressed images (92.7\%). These aggregate rates concealed strong crop differences (Table~\ref{tab:coverage}). Ash gourd retained 996 of 997 images, compared with 580 of 798 for tomato. Tomato alone accounted for 218 of the 362 exclusions (60.2\%). The leaf-miner category retained 107 of 207 images (51.7\%). All descriptor-model estimates therefore apply to a sample depleted in this phenotype. Of the 362 excluded masks, 349 fell below the minimum area and 13 exceeded the maximum.
\showpanel{\cohorttable}

\subsection{Descriptor shifts depended on crop identity}
The direction of a descriptor shift could change across crops (Figure~\ref{fig:traits}). For contour area the signed rank contrast was $-0.64$ in cucumber and $+0.40$ in eggplant. For texture contrast, ash gourd gave $+0.57$ and tomato $-0.34$. Edge yellow ratio showed a larger contrast in bottle gourd ($+0.49$) than in bitter gourd ($+0.03$). A pooled rule that treats an increase in one measurement as evidence of stress will therefore meet different class distributions in different crops. These are within-crop comparisons of recorded image values. The stress category combines several conditions whose frequencies differ across crops, and camera distance, leaf orientation and the background inside the texture window also influence what is measured. An area contrast, for example, can arise from scale as well as from leaf size. The contrasts motivate crop-specific classifiers; they do not attribute a growth response to stress.
\showpanel{\traitpanel}

\subsection{Effect of topology and refinement}
Figure~\ref{fig:depth}A compares the 20 pooled configurations. All five graphs had the same macro-F1 at depth zero (0.601). By depth three the hybrid graph reached 0.697 and the family graph 0.691, the nearest-neighbor graph 0.674, and the correlation and regional graphs approximately 0.622. These maxima come from the sweep used to select the configuration.
\showpanel{\depthpanel}

\subsection{What the refinement gain consists of}
Tuning the SVM cost did not remove the gain (Table~\ref{tab:mechanism}). With nested cost selection, descriptor states alone reached 0.596 pooled macro-F1, and the hybrid graph at depth three again selected $C=1$ in every fold (0.697). The refinement gain was therefore not an artifact of a fixed cost. It was, however, much smaller in balanced accuracy (0.698 to 0.722) than in macro-F1, because the two representations made different errors. Descriptor states alone recognized 73.7\% of healthy but only 65.9\% of stressed leaves, whereas the depth-three kernel recognized 58.2\% and 86.1\%. Since stressed leaves are the majority, macro-F1 rewards the second trade-off more.

Most deep labels never recurred. On the hybrid graph, 90.7\% of held-out node labels matched the training vocabulary at depth one, 33.3\% at depth two and 4.4\% at depth three, and 95.9\% of depth-three training labels occurred in a single training image (S23). The family graph behaved identically at depth three. Sparse graphs retained matches (correlation 91.8\%, regional 97.1\% at depth three) but gained little. The useful shared information therefore came from depths one and two, while depth three mainly added coordinates unique to each training image.

The private-pattern control, which reproduces that uniqueness without any shared neighborhood, raised pooled macro-F1 from 0.596 to 0.629, about a third of the 0.096 refinement gain. A kernel counting matched pairs of descriptor states, with no graph, reached 0.652, and one counting matched triples reached 0.690, within 0.007 of the hybrid WL kernel. The same ordering held within crops and with predicted crop (Table~\ref{tab:mechanism}): the graph-free triple kernel scored 0.699 and 0.683, compared with 0.685 and 0.682 for nested WL. Extending the WL grid to depth zero and tuned cost lowered the nested scores slightly (0.679 and 0.681), consistent with selection noise in a larger grid.
\showpanel{\mechtable}

\subsection{Random graphs with the hybrid edge count}
Figure~\ref{fig:depth}B compares the hybrid graph with 200 random graphs (Table~\ref{tab:mechanism}). Permuted descriptor positions gave a mean five-fold macro-F1 of 0.695 (range 0.680--0.712), degree-preserving rewiring 0.695 (0.684--0.702) and random graphs with the same edge count 0.698 (0.689--0.708). The hybrid graph (0.697) was matched or exceeded by 41 of 100 permuted, 18 of 50 rewired and 29 of 50 random graphs (one-sided empirical $p=0.42$, $0.37$ and $0.59$; S22b). Its balanced accuracy (0.722) was modestly higher than the null means (0.701--0.705), so the proposed edges shifted the balance of class errors slightly but did not improve macro-F1. Any graph with the hybrid graph's density supported the same gain. The vocabulary of reassigned graphs was also larger, because the number of distinct combinations depends on which measurements are joined (S22).

\subsection{Graph construction affected the four graph models differently}
WL ranged from 0.621 to 0.697 across the five graphs at fixed depth (Table~\ref{tab:graphstudy}). GIN scored 0.664 without inter-node edges and 0.672 on the hybrid graph. GCN achieved its highest mean without inter-node edges (0.667). GAT improved from 0.601 without inter-node edges to 0.683 on the nearest-neighbor graph, where it exceeded WL (0.674). All three networks exceeded WL on the correlation and regional graphs, and WL led on the hybrid and family graphs. Learned identity embeddings and pooled readout give the networks access to image-wide combinations even without messages between nodes, which may explain why their node-only controls remained competitive. Averaging over the five shared graphs and folds, a WL fit took 9.6~s on CPU and GIN, GCN and GAT fits 59--76~s on a GPU; these times combine architecture, training procedure and device and are not a like-for-like speed comparison.
\showpanel{\graphstudytable}

\subsection{A held-out leaf through the kernel}
A single held-out example illustrates what the kernel accumulates. We selected the first healthy cucumber image in lexicographic filename order from the first test fold, independently of its prediction. In image 15 the contour occupied 4,043,665 square pixels, in the upper training tertile; mean hue was in the upper bin, green ratio in the lower bin and texture contrast in the middle bin. All 38 initial states had entries in the training vocabulary. At the first refinement 35 of the 38 nodes matched a training neighborhood, at the second 21 and at the third none. The SVM returned a decision margin of $-0.029$, on the healthy side of zero, based entirely on the initial and two shallower levels. S14 gives the measured values and states.

\subsection{Comparison with descriptor classifiers and image encoders}
When crop identity was supplied, tuned Extra Trees had the highest mean crop macro-F1 (0.742), followed by MobileNetV3 (0.734), tuned XGBoost (0.731) and EfficientNet-B0 (0.729) (Table~\ref{tab:tasks}). Selected WL reached 0.716 and nested WL 0.685. Individual crops showed a different ordering (Figure~\ref{fig:crops}): nested WL scored 0.841 for cucumber but 0.541 for ridge gourd, where EfficientNet-B0 reached 0.720. The comparison cannot isolate a species effect from class composition or acquisition differences.
\showpanel{\croppanel}

In the pooled task histogram gradient boosting achieved 0.729, above selected WL (0.697), nested WL (0.691) and EfficientNet-B0 (0.705). With predicted crop, frozen EfficientNet-B0 achieved the highest mean macro-F1 (0.739), followed by MobileNetV3-Large (0.734) and tuned Extra Trees (0.729); nested WL reached 0.683 (Table~\ref{tab:routed}). The always-stressed rule achieved 83.8\% accuracy with zero healthy recall. The default random forest exceeded EfficientNet-B0 in accuracy but detected only 31.3\% of healthy leaves, compared with 62.0\% for EfficientNet-B0. Nested WL detected 53.9\% of healthy leaves and 86.5\% of stressed leaves. Supplying the true crop changed nested WL from 0.683 to 0.686 and tuned Extra Trees from 0.729 to 0.736, so routing errors were not the main source of the deficit. The descriptor router was correct for 91.9\% of held-out images and the EfficientNet router for 99.2\%.
\showpanel{\tasktable}
\showpanel{\routedtable}

\subsection{Matched comparisons}
Within crops, selected WL led the three default GNNs on 31--32 of 40 folds (Table~\ref{tab:paired}), but this comparison favors WL because only its configuration was chosen on the study data. Under nested selection for both, WL exceeded GIN by 0.052 macro-F1; the 40-fold test gave $p=0.004$ and the eight-crop test $p=0.148$. Nested WL trailed the graph-free triple kernel by 0.014 (17 of 40 folds won; fold-level $p=0.112$, crop-level $p=0.195$), tuned Extra Trees by 0.057 (10 of 40; $p<0.0001$ and $p=0.039$) and EfficientNet-B0 by 0.044 (13 of 40; $p=0.020$ and $p=0.312$). Default and tuned XGBoost scored 0.005 and 0.015 above selected WL. For pooled default XGBoost, edge mean hue had the largest mean absolute TreeSHAP contribution (0.467 log-odds), followed by hue variability (0.450); shape contributed 25.6\% of total absolute attribution and Lab color 19.4\% (S18--S19).
\showpanel{\pairedtable}

\subsection{Where a positive stress screen was lost}
The pooled depth-three hybrid returned positive calls on 3,165 of 3,674 accepted stressed images and negative calls on 415 of 713 accepted healthy images, leaving 509 missed stress cases and 298 false positives. Potassium deficiency contributed 115 misses, nitrogen deficiency 96 and mite damage 85, together 58.2\% of misses. Their denominators differed: mite damage was missed on 85 of 257 accepted images, nitrogen deficiency on 96 of 673.

Table~\ref{tab:errors} follows individual conditions from submission to screen. Of 207 leaf-miner photographs, 100 were excluded before classification and another 19 accepted images received a healthy prediction, so the descriptor pipeline returned a positive screen for 88 of 207 (42.5\%) even though it recognized 88 of the 107 retained examples (82.2\%). EfficientNet-B0, trained on the same accepted photographs but scoring every submitted photograph, returned a positive screen for 188 of 207 leaf-miner photographs (90.8\%), including 95 of the 100 that segmentation had rejected; among the 107 accepted leaf-miner photographs it flagged 93, compared with 88 for WL. For pumpkin leaf-eating insect damage it flagged 48 of 80 photographs (60.0\%), compared with 35 (43.8\%) for the descriptor pipeline. The rejected photographs were not intrinsically harder to classify: on the 362 rejected photographs alone, EfficientNet-B0 reached a mean fold macro-F1 of 0.766 (80.7\% of stressed and 86.3\% of healthy leaves recognized), and its macro-F1 over all submitted photographs (0.711) was similar to that on the accepted cohort (0.705, matching the archived pooled result). Adding the rejected training photographs changed these values by less than 0.01 (S25). MobileNetV3-Large behaved similarly (172 of 207 leaf-miner photographs). These positive calls indicate any recorded stress; recognizing the cause would require a separate prediction.
\showpanel{\errtable}

\subsection{Acquisition changes affected both coverage and recognition}
Repeating the robustness study on all five test folds changed its conclusion (Table~\ref{tab:robust}, Figure~\ref{fig:robust}). Fold zero reproduced the original single-fold result for WL ($-0.053$), but across folds the mean change was $-0.088\pm0.032$ for WL, $-0.083\pm0.032$ for nearest neighbors, $-0.089\pm0.005$ for the MLP and $-0.095\pm0.034$ for tuned Extra Trees; tuned XGBoost lost $0.113\pm0.015$ and the default random forest $0.137\pm0.033$. WL had the smallest decline among the twelve descriptor models in two folds and ranked tenth in another, and its paired difference from tuned Extra Trees and from nearest neighbors changed sign across folds. WL therefore belonged to the more robust group of descriptor models without being reliably the most robust. Frozen encoders were nearly unaffected: $-0.003\pm0.009$ for EfficientNet-B0 and $-0.014\pm0.016$ for MobileNetV3. On the 1,083 photographs with predictions from every model under every change (215--217 per fold), the ordering was the same (WL $-0.090$, nearest neighbors $-0.079$, tuned Extra Trees $-0.100$, EfficientNet-B0 $-0.003$).

The warm cast was the most damaging change for WL ($-0.154\pm0.047$) and for tuned Extra Trees ($-0.252$). Changes also removed photographs before classification: of 1,201 photographs per change, segmentation rejected 80 after brightening, 57 after the warm cast, 23 after darkening, 14 after noise and 2 after JPEG compression, while blur and halving rejected none. Brightening thus cost coverage while leaving macro-F1 on the retained photographs comparatively intact. In the original single-fold study the GNNs lost 0.068 (GIN), 0.089 (GCN) and 0.095 (GAT) on average, compared with 0.053 for WL on the same fold (Figure~\ref{fig:roundthree}B); given the fold-to-fold variation above, that single-fold ranking should not be generalized.
\showpanel{\robusttable}
\showpanel{\robustpanel}

\subsection{Training data and single-image processing time}
At 20\% of the within-crop training data, selected WL attained macro-F1 of 0.593, the radial-basis SVM 0.632 and MobileNetV3 0.648 (Figure~\ref{fig:cost}). With all training examples, WL reached 0.716 and the radial-basis SVM 0.679. At 10\% of the data selected WL achieved 0.555, compared with 0.628 for MobileNetV3. Default XGBoost exceeded selected WL at all six fractions, and selected WL exceeded each default GNN at all six (Figure~\ref{fig:roundthree}A). Because the WL settings were selected on the study data, these curves favor WL. Frozen encoders also benefit from external pretraining; the curves measure demand for labeled target images given that pretraining. For routed inference the estimated times were 1,077~ms per image for selected WL, 1,074~ms for tuned Extra Trees and 143~ms for EfficientNet-B0. Reading, segmenting and measuring the full-resolution image accounted for approximately 1,064~ms of the descriptor pipeline, so a faster final classifier would address little of that cost.
\showpanel{\costpanel}
\showpanel{\roundthreepanel}

\section{Discussion}
\subsection{What the graph kernel contributed}
Comparing descriptor states jointly, rather than one at a time, produced the largest gain we observed: from 0.596 for states alone to 0.690 for a graph-free kernel on triples of states and 0.697 for the hybrid WL kernel. The graph itself added little beyond that. A third of the refinement gain was reproduced by a control in which no neighborhood is ever shared, because deep WL patterns were almost all unique to one training image; this mechanism resembles the diagonal dominance described for other sparse kernels \citep{Yanardag2015,Greene2006}. Random graphs with the same number of edges performed as well as the hybrid graph, whether they kept its shape, its degree sequence or neither. Under nested selection WL did not outperform the graph-free kernel within crops or with routing. The measured benefit therefore belongs to joint matching of discretized descriptors, not to a particular proposed link between traits.

Several reasons may explain why the specific edges mattered so little. Color descriptors in different spaces encode closely related signals, so replacing one neighbor with another correlated measurement may preserve information even when an edge changes its name. Sparse graphs kept their labels matchable but combined few descriptors, whereas dense graphs combined many descriptors but rapidly produced unique labels; the best graphs sat between these regimes at depths one and two. The comparison of graph models points the same way: GCN and GIN were competitive without any inter-node edges, and GAT preferred a different graph from WL. Tree ensembles, which form descriptor interactions adaptively, reached the highest descriptor scores, in line with benchmarks of tabular data \citep{Grinsztajn2022,ShwartzZiv2022}. For future graph designs on measurement vectors, a graph-free interaction kernel and random graphs with matched edge counts are inexpensive controls that should be reported alongside any proposed topology.

\subsection{Interpretability and the limits of condition recognition}
For a given image, the WL representation shows which quantile each measurement fell into and which neighborhoods matched training patterns, and a reviewer can trace a match to descriptors such as core hue and edge yellow ratio. Such a trace explains a computational decision. Quantile states depend on the training cohort, and an association with the stress label cannot establish a nutritional deficiency or justify treatment. Because most deep patterns are unique, a trace is informative mainly at the shallow depths where patterns recur. Localized lesions also occupy a small fraction of a leaf and may contribute little to whole-leaf summaries; the image encoders are consistent with the value of retaining spatial detail, although resolution, pretraining, representation and classifier all differ between systems.

The binary label also limits what a correct match establishes. In the project's separate 57-class reanalysis of the same accepted feature table, the highest topology/depth mean was 0.428 macro-F1 (hybrid graph, depth two), compared with 0.550 for a flat random forest and 0.472 for a linear SVM (S13). The binary endpoint gives full credit for calling two entirely different injuries stressed: the 78 positive calls among 78 accepted powdery-mildew images show that those images appeared abnormal, not that mildew was diagnosed.

\subsection{Phenotyping implications}
Segmentation decided which leaves could be screened at all. The excess-green index rejected about half of the leaf-miner photographs and more than a quarter of tomato photographs. A frozen encoder that needs no mask, trained on exactly the same photographs, recovered 95 of the 100 rejected leaf-miner screens and classified rejected photographs as well as accepted ones. The coverage loss was therefore a property of the measurement pipeline, not of the photographs. A measurement pipeline that reports only the accuracy on accepted photographs therefore overstates its usefulness for exactly the phenotypes that change leaf color or integrity. A screening study should report both denominators, accepted and submitted, and keep a log of rejected images; obtaining expert outlines for a stratified sample of accepted and rejected leaves would distinguish color loss, fragmented margins and background artifacts as causes of rejection.

The operating point matters as much as the ranking. A model that labels almost every leaf stressed has high sensitivity in this collection and imposes a substantial review burden on healthy samples, and even the best routed system classified about 38\% of healthy images as stressed. The observed prevalence of stress is specific to this collection, so predictive values and review workload would change elsewhere.

Acquisition changes act on two components at once: a color cast can move a leaf across the mask-area threshold, and it can move a retained leaf across a training quantile. Quantization dampens only changes that stay within a bin. Across five folds, all descriptor pipelines lost roughly 0.08--0.15 macro-F1 under these changes, the warm cast was the most damaging, and the ranking among descriptor models varied by fold, whereas frozen encoders were almost unaffected. A single test fold was not enough to rank robustness; we recommend repeating perturbation studies across folds.

\subsection{Limitations}
The folds establish performance within one collection; they do not establish transfer to new plants, sites, devices or backgrounds. OLID-I uses a single primary camera and a controlled black background, and the image-level split leaves possible dependence among photographs of the same plant or session unresolved, which may inflate all scores. Synthetic changes probe sensitivity to specified transformations, whereas field photographs also vary viewpoint, occlusion and symptom mixtures. Pixel areas and the texture window carry camera distance and leaf outline, which may act as crop cues rather than stress cues. The selected WL configuration was chosen on the study data, and the learning curves and GNN comparisons that use it favor WL. GCN and GAT were not tuned. The most informative external test would retain field identifiers, score whole submitted cohorts and prespecify the representation and decision rule before evaluation.

\subsection{Conclusion}
On 38 handcrafted leaf measurements, the WL kernel's gain over individual descriptor states came mainly from comparing combinations of discretized descriptors, partly from the scale effect of patterns unique to single training images, and little from the specific graph: a graph-free kernel on triples of states performed as well under nested selection, and 200 random graphs with the same edge count performed as well as the proposed hybrid graph. WL was competitive with but not better than tuned tree ensembles, belonged to the more robust group of descriptor models without being reliably the most robust, and was outperformed by frozen image encoders in routed accuracy, robustness to acquisition changes and processing time. Segmentation removed about half of the leaf-miner photographs before classification, and an encoder that needs no mask recovered almost all of those screens. Studies of image-based stress screening should report how many photographs reach the classifier, compare proposed graphs with graph-free and random-graph controls, and select every compared model by the same nested procedure.

\section*{Data availability statement}
The OLID-I photographs and annotations are deposited by Orka and colleagues at \url{https://doi.org/10.5281/zenodo.8105154}. The project code and derived measurements are maintained at \url{https://github.com/BharathBala21/WL_Graph_Representation}; \textit{an archived release with a persistent identifier will be cited here before submission (author action)}. Supplementary Data Sheet S1 records the 4,387 image identifiers and fixed pooled/routed folds. S2 gives all routed fold scores, S3 the pooled model comparison, S4 the crop-specific comparison, S5 single-fold robustness estimates, S6 learning curves, S7 the descriptor dictionary, S8 the topology sweep and S9 the single-fold common-image analysis. S10 contains the original five-permutation controls, S11 the original-hybrid held-out calls, S12 the crop-specific contrasts, S13 the source summaries of the 57-class comparison and S14 the worked example. S15--S17 provide known-crop folds, the graph study and paired comparisons, and S18--S19 XGBoost attributions. S20 gives the cost and graph-free kernel grids, S21 nested selections for the pooled and known-crop tasks, S22 the random-graph controls, S23 vocabulary matching by depth and S24 routed nested results. S25--S26 report the image encoders on every submitted photograph, and S27--S28 the five-fold robustness study.

\section*{Author contributions}
\textit{Author confirmation required: individual CRediT roles, byline order and corresponding author.}
\section*{Funding}
\textit{Author confirmation required: funding sources and award numbers, or a declaration of no specific funding.}
\section*{Conflict of interest}
\textit{Author confirmation required: declarations from all authors.}
\section*{Generative AI statement}
The authors used generative AI tools during this work. OpenAI Codex (GPT-6.2 Astra, OpenAI) was used in November 2026 to draft an earlier version of the manuscript, prepare LaTeX and plotting code, check reported metrics against project artifacts and edit for clarity. Claude Code (Claude Opus 5.5, Anthropic) was used in November 2026 to review the manuscript, write and run the analysis code for the revision experiments (cost tuning, graph-free and private-pattern kernels, random-graph controls, encoder coverage and five-fold robustness), and revise the manuscript text. The authors reviewed and verified all code, results and text and take full responsibility for the content of this article.
\section*{Acknowledgments}
The authors acknowledge the OLID-I dataset creators for releasing the images and annotations.
